\chapter{Wstęp}

\section{Motywacja}

Rynek aplikacji mobilnych stał się dominującym kanałem konsumpcji treści i usług cyfrowych. Dane o penetracji dostępu mobilnego (ok.\ 71\% globalnej populacji posiada subskrypcję mobilną do 2023 r.) \cite{CiscoAIRWhitepaper2023} oraz statystyki udziału ruchu WWW (ponad 50\% globalnych \emph{page views} pochodzi z urządzeń mobilnych) \cite{StatCounterPlatformShare2025} pokazują, że to właśnie środowisko mobilne jest dziś kluczowym miejscem, w którym kształtują się wzorce komunikacji sieciowej i materializują ryzyka dla prywatności.

W ekosystemie Androida współistnieje wiele bibliotek i zestawów SDK firm trzecich — od analityki i reklam, przez personalizację, po rozwiązania antyfraudowe. Każda z tych warstw wnosi własne punkty końcowe, protokoły i mechanizmy telemetrii, często działające poza świadomością użytkownika, a niekiedy również poza pełną kontrolą zespołów deweloperskich aplikacji. W efekcie rośnie złożoność oraz nieprzejrzystość przepływów danych.

Konsekwencją tego stanu rzeczy są konkretne wyzwania badawcze:
\begin{enumerate}[label=\alph*)]
  \item uzyskanie przejrzystości, \emph{jakie} protokoły, domeny i dostawcy są realnie wykorzystywani przez aplikacje,
  \item ocena ryzyka ekspozycji metadanych (np.\ hosty, czasy, rozmiary transferów) mimo powszechnego szyfrowania treści,
  \item porównywalność zachowań komunikacyjnych między kategoriami aplikacji (komunikatory, multimedia, finanse, gry),
  \item budowa powtarzalnych narzędzi i procedur, które nadążą za ewolucją protokołów (HTTP/2, HTTP/3/QUIC) oraz mechanizmami obronnymi (np.\ \emph{certificate pinning}).
\end{enumerate}

Wobec powyższego potrzebne jest uporządkowane, możliwie obiektywne podejście do \emph{systematycznej} obserwacji i oceny ruchu sieciowego aplikacji: środowisko pomiarowe, które pozwala porównywać różne aplikacje i kategorie, zdefiniowany zestaw metryk (zarówno ilościowych, jak i jakościowych) oraz protokoły postępowania uwzględniające ograniczenia współczesnych technologii transportowych i zabezpieczeń. Tak rozumiana ramka badawcza stanowi bezpośrednią motywację dla niniejszej pracy.


Drugi wektor motywacyjny stanowi szybka adopcja nowych protokołów transportowych i warstw aplikacyjnych. QUIC (standaryzowany jako HTTP/3 \cite{RFC9114}) przenosi część logiki transportu do przestrzeni użytkownika (UDP + TLS 1.3 \cite{RFC8446}), co zmienia charakterystykę obserwowalnych cech przepływów: klasyczne estymatory opóźnień i retransmisji z poziomu TCP stają się mniej bezpośrednie, a analiza pakietów jest utrudniona przez szyfrowanie większej liczby pól. W konsekwencji wcześniejsze proporcje protokołów publikowane w literaturze mogą być nieaktualne dla obecnych wersji popularnych aplikacji.

Kolejnym aspektem jest prywatność i otoczenie regulacyjne. RODO (GDPR), a także DMA/DSA w UE akcentują minimalizację i przejrzystość przetwarzania. Architektura „multi-endpoint” (API własne + analityka + reklamy + antyfraud + CDN/edge) prowadzi do gęstych, równoległych zestawów połączeń tuż po interakcji użytkownika. Nawet przy pełnym szyfrowaniu treści korelacja metadanych (częstotliwość, rozmiary odpowiedzi, domeny) może ujawniać wzorce użycia funkcji (pattern-of-life). Transparentność takich wzorców jest ograniczona zarówno dla użytkownika, jak i audytorów. Narzędzia systemowe dają uogólniony obraz, a techniki inwazyjne (hooking, dynamic instrumentation) bywają blokowane przez mechanizmy anty‑tampering.

Istnieje także motywacja praktyczna stworzenia lekkiego, powtarzalnego środowiska badawczego balansującego między:
\begin{itemize}
  \item \textbf{Zasięgiem obserwacji} (kompletność metadanych HTTP[S]),
  \item \textbf{Inwazyjnością} (brak konieczności rootowania, minimalne modyfikacje),
  \item \textbf{Zgodnością prawną i etyczną} (anonimizacja identyfikatorów, brak utrwalania treści wrażliwych).
\end{itemize}
Istniejące narzędzia (np. \emph{mitmproxy}) zapewniają bazowy przechwyt, ale wymagają rozszerzeń o: 
\begin{itemize}
  \item \textbf{Normalizacja} – ujednolicenie formatu danych.
  \item \textbf{Klasyfikacja} – kategoryzacja ruchu sieciowego.
  \item \textbf{Agregacja} – zbieranie i analiza metryk.
  \item \textbf{Automatyzacja} – skrypty do interakcji z aplikacjami.
  \item \textbf{Anonimizacja} – ochrona danych osobowych.
\end{itemize}

Z naukowego punktu widzenia identyfikowane luki obejmują:
\begin{description}
  \item[L1] Brak bieżących, empirycznych pomiarów dystrybucji HTTP/1.1 vs HTTP/2 vs QUIC w popularnych aplikacjach.
  \item[L2] Ograniczone porównania między kategoriami aplikacji a zachowaniami komunikacyjnymi.
  \item[L3] Niedostatek jawnych metryk ryzyka prywatności bazujących wyłącznie na metadanych.
  \item[L4] Fragmentaryczne dane o wpływie certificate pinning na pokrycie przechwytu.
\end{description}

Interesariusze:
\begin{itemize}
  \item \textbf{Użytkownik} – zyskuje pośrednio poprzez privacy-by-design.
  \item \textbf{Deweloper} – możliwość audytu nadmiarowych zależności.
  \item \textbf{Audytor bezpieczeństwa} – metryki priorytetyzacji ryzyka.
  \item \textbf{Społeczność naukowa} – powtarzalny pipeline do dalszych badań.
\end{itemize}

\section{Stan wiedzy}

\subsection{Analiza ruchu sieciowego}
Badania pokazują, że ruch generowany przez aplikacje mobilne niesie znaczące wyzwania dla prywatności – nawet konfiguracje TLS nie chronią całkowicie, bo analizę można przeprowadzać na metadanych lub kanałach bocznych. W przeglądzie literatury poświęconym analizie ruchu mobilnego autorzy klasyfikują cele, miejsca przechwytywania oraz metody wykorzystania danych sieciowych do profilowania użytkowników \cite{Conti2017TrafficAnalysis}. Inne prace demonstrują, jak sieć może zostać wykorzystana do identyfikacji konkretnej aplikacji, nawet przy szyfrowaniu, wykorzystując uczenie maszynowe na podstawie cech pakietów \cite{Taylor2017AppFingerprinting}. Przykłady wykorzystania narzędzi takich jak \textit{mitmproxy} i \textit{tcpdump} pokazują, jak można wychwycić i analizować zarówno HTTP jak i inne kanały sieciowe \cite{mitmproxy,tcpdump}.

\subsection{Inżynieria wsteczna i testy penetracyjne}
Inżynieria wsteczna aplikacji mobilnych, ułatwiająca badanie mechanizmów takich jak SSL pinning, jest szeroko stosowana. Wczesne prace wskazują na powszechne błędy konfiguracji SSL w aplikacjach Androidowych — ponad połowa analizowanych aplikacji korzystających z SSL miała błędne wdrożenia lub brak obsługi poprawnego pinningu \cite{Fahl2012WhyEve}. W celu przeprowadzania dynamicznych testów bezpieczeństwa wykorzystuje się narzędzia takie jak \textit{Frida}, \textit{Drozer} czy \textit{MobSF}, które wspierają zarówno analizę dynamiczną, jak i automatyzację \cite{MobSF,Frida}.

\subsection{Prywatność, uprawnienia i trackery}
Coraz więcej badań zwraca uwagę na skalę gromadzenia danych przez aplikacje mobilne. Systematyczne opracowania przeglądowe pokazują spektrum metod oceny prywatności (statycznych, dynamicznych i hybrydowych) oraz wskazują luki w zgodności z regulacjami \cite{DelAlamo2021PrivacyMapping}. Narzędzia takie jak \textit{Exodus Privacy} umożliwiają analizę APK pod kątem wbudowanych trackerów i żądanych uprawnień \cite{ExodusPrivacy}. Wyniki badań empirycznych pokazują, że aplikacje często żądają nadmiarowych uprawnień lub integrują zewnętrzne trackery reklamowe \cite{Razaghpanah2018AppsPrivacy}.

\subsection{Zużycie zasobów i kontekst behawioralny}
Nadmierne użycie zasobów (np. bateria, CPU) może wskazywać na obecność komponentów śledzących lub brak optymalizacji. Do pomiarów wykorzystuje się m.in. \textit{Battery Historian} oraz \textit{Android Profiler} \cite{BatteryHistorian,AndroidProfiler}. Chociaż ten aspekt jest rzadziej eksplorowany w literaturze naukowej, podkreśla się jego znaczenie w identyfikacji ukrytych mechanizmów aplikacji \cite{Pathak2012BatteryDrain}.

\subsection{Podsumowanie}
Przegląd literatury wykazuje, że istnieje bogata baza wiedzy w obszarach analizy ruchu sieciowego, inżynierii wstecznej, testów penetracyjnych oraz badania prywatności aplikacji mobilnych. Jednocześnie wciąż brakuje holistycznego podejścia, które łączyłoby te wszystkie elementy — co stanowi niszę, którą niniejsza praca ma ambicję wypełnić.


\section{Cel i zakres pracy}
\cel{Celem pracy jest opracowanie i ewaluacja środowiska badawczego umożliwiającego systematyczną analizę bezpieczeństwa i prywatności aplikacji mobilnych na system Android. Praca koncentruje się na przechwytywaniu, klasyfikacji oraz ocenie ruchu sieciowego, a także na badaniu mechanizmów uwierzytelniania, uprawnień i zużycia zasobów.}

Zakres pracy obejmuje:
\begin{itemize}
  \item \textbf{Środowisko przechwytywania} – konfiguracja narzędzi takich jak \textit{mitmproxy}, \textit{Wireshark}, \textit{tcpdump}, a także przygotowanie środowiska laboratoryjnego z wykorzystaniem emulatorów oraz zrootowanych urządzeń mobilnych.
  \item \textbf{Analiza ruchu sieciowego} – identyfikacja używanych protokołów (HTTP/1.1, HTTP/2, QUIC, TLS), badanie bezpieczeństwa transmisji danych (szyfrowanie, SSL pinning) oraz klasyfikacja przesyłanych informacji.
  \item \textbf{Reverse engineering i testy penetracyjne} – dekompilacja aplikacji z wykorzystaniem \textit{APKTool}, \textit{Jadx}, \textit{Ghidra} oraz testy przy użyciu narzędzi takich jak \textit{MobSF}, \textit{Frida}, \textit{Metasploit}.
  \item \textbf{Analiza uprawnień i prywatności} – badanie przydzielanych uprawnień, wykrywanie trackerów i komponentów reklamowych (\textit{Exodus Privacy}, \textit{App Inspector}) oraz ocena zarządzania danymi użytkowników.
  \item \textbf{Analiza zużycia zasobów} – pomiary obciążenia baterii i CPU aplikacji z użyciem \textit{Battery Historian} oraz \textit{Android Profiler}.
  \item \textbf{Eksperymenty} – testy na wybranych aplikacjach z różnych kategorii (m.in. komunikatory, aplikacje społecznościowe, bankowe, rządowe), analiza różnic w implementacji mechanizmów bezpieczeństwa.
  \item \textbf{Badanie wektorów ataków} – symulacja ataków typu \textit{Man-in-the-Middle}, analiza przechowywania danych lokalnych, scenariusze działania offline/online, testy odporności na automatyczne skrypty i boty.
  % \item \textbf{Analiza logów i ścieżki audytu} – sprawdzenie, jakie dane aplikacje rejestrują w logach oraz ocena ryzyka wycieku informacji.
\end{itemize}
